
\documentclass{article}
\usepackage{colortbl}
\usepackage{makecell}
\usepackage{multirow}
\usepackage{supertabular}

\begin{document}

\newcounter{utterance}

\centering \large Interaction Transcript for game `clean\_up', experiment `1\_medium\_7obj\_en', episode 1 with qwen3{-}next{-}80b{-}a3b{-}instruct{-}t1.0.
\vspace{24pt}

{ \footnotesize  \setcounter{utterance}{1}
\setlength{\tabcolsep}{0pt}
\begin{supertabular}{c@{$\;$}|p{.15\linewidth}@{}p{.15\linewidth}p{.15\linewidth}p{.15\linewidth}p{.15\linewidth}p{.15\linewidth}}
    \# & $\;$A & \multicolumn{4}{c}{Game Master} & $\;\:$B\\
    \hline

    \theutterance \stepcounter{utterance}  
    & & \multicolumn{4}{p{0.6\linewidth}}{
        \cellcolor[rgb]{0.9,0.9,0.9}{
            \makecell[{{p{\linewidth}}}]{
                \texttt{\tiny{[P1$\langle$GM]}}
                \texttt{Sono il tuo game master, e stai giocando un gioco collaborativo con la seguente griglia come tavolo da gioco:} \\
\\ 
\texttt{```} \\
\texttt{ 1234567} \\
\texttt{╔═══╤═╤═╗} \\
\texttt{║◌◌◌│◌│◌║ 1} \\
\texttt{║◌◌M│◌├─╢ 2} \\
\texttt{║AT◌├─┘◌║ 3} \\
\texttt{║◌◌P│◌◌◌║ 4} \\
\texttt{╟─┬─┼───╢ 5} \\
\texttt{║◌│◌│S◌O║ 6} \\
\texttt{║◌│◌│◌D◌║ 7} \\
\texttt{╚═╧═╧═══╝} \\
\\ 
\texttt{```} \\
\\ 
\texttt{* Il bordo superiore mostra le coordinate x che aumentano verso destra, e il bordo destro le coordinate y che aumentano verso il basso.} \\
\texttt{* The following objects are randomly placed on your grid: 'P', 'O', 'T', 'S', 'D', 'A', 'M'.} \\
\\ 
\\ 
\texttt{L'altro giocatore vede una variante del tavolo da gioco, dove gli oggetti sono posizionati in diverse posizioni casuali. Non puoi vedere il tavolo dell'altro giocatore, e loro non possono vedere il tuo.} \\
\\ 
\\ 
\texttt{**Obiettivo:**} \\
\\ 
\texttt{Entrambi i giocatori devono spostare gli oggetti sui rispettivi sfondi in modo che oggetti identici finiscano alle stesse coordinate. Devi comunicare con l'altro giocatore per concordare una configurazione obiettivo comune.} \\
\\ 
\\ 
\texttt{**Regole:**} \\
\\ 
\texttt{* In ogni turno, puoi inviare esattamente uno dei seguenti due comandi:} \\
\texttt{ 1. `COMUNICA: <MESSAGE>`: per inviare un messaggio (tutto fino alla prossima interruzione di riga) all'altro giocatore. Lo inoltrerò al tuo partner.} \\
\texttt{2. `SPOSTA: <OBJECT>, (<X>, <Y>)`: per spostare un oggetto in una nuova posizione, dove <X> è la colonna e <Y> è la riga. Ti farò sapere se la tua mossa è andata a buon fine o no.} \\
\texttt{* Se non rispetti il formato, o invii diversi comandi contemporaneamente, devo penalizzarti.} \\
\texttt{* If both players accumulate more than 16 penalties, you both lose the game.} \\
\texttt{* È vitale che comunichi con l'altro giocatore riguardo al tuo stato obiettivo! L'*unico* modo in cui puoi trasmettere la tua strategia all'altro giocatore è usando il comando `COMUNICA: <MESSAGE>`!} \\
\\ 
\\ 
\texttt{**Spostamento Oggetti**} \\
\\ 
\texttt{* Puoi spostare oggetti solo in celle entro i limiti della griglia. La cella di destinazione deve essere vuota, cioè deve contenere solo il simbolo '◌'.} \\
\texttt{* Se provi a spostare un oggetto in un posto che non è vuoto, o provi a spostarlo fuori dalla griglia, devo penalizzarti. Hai un altro tentativo.} \\
\texttt{* Prima di fare una mossa, controlla due volte che il posto di destinazione sia vuoto, e non contenga alcuna lettera, cornice, o linea!} \\
\\ 
\\ 
\texttt{**Fine del Gioco**} \\
\\ 
\texttt{Se pensi di aver raggiunto l'obiettivo di allineare tutti gli oggetti, puoi chiedere all'altro giocatore di finire il gioco inviando COMUNICA: finito?. Se l'altro giocatore ti chiede di finire il gioco, e rispondi con COMUNICA: finito!, the game will end.} \\
\\ 
\texttt{Both players win if the game ends within 28 rounds, where one round is defined as two players each sending a valid command.} \\
\\ 
\texttt{**Punteggio:**} \\
\\ 
\texttt{Più vicini sono gli oggetti identici in entrambi i tavoli da gioco, più punti ottieni. Le penalità riducono i tuoi punti. Riesci a battere il record?} \\
\\ 
\\ 
\texttt{Per favore invia un messaggio all'altro giocatore per iniziare il gioco!} \\
            }
        }
    }
    & \\ \\

    \theutterance \stepcounter{utterance}  
    & \multicolumn{4}{p{0.6\linewidth}}{
        \cellcolor[rgb]{0.9,0.9,0.9}{
            \makecell[{{p{\linewidth}}}]{
                \texttt{\tiny{[P1$\rangle$GM]}}
                \texttt{COMUNICA: Ciao! Il mio tavolo ha questi oggetti: P(4,4), O(7,6), T(1,3), S(5,6), D(7,7), A(1,3), M(3,2). Nota che A e T sono entrambi in riga 3, colonne 1 e 3. Ho bisogno di allineare tutti gli oggetti con te. Per favore, descrivi la posizione dei tuoi oggetti e proporremo una configurazione obiettivo comune. Potremmo mirare a sistemarli in una riga o colonna libera, ad esempio la riga 5 o 6 che ha più spazio vuoto. Aspetto la tua risposta!} \\
            }
        }
    }
    & & \\ \\

    \theutterance \stepcounter{utterance}  
    & & & \multicolumn{2}{p{0.3\linewidth}}{
        \cellcolor[rgb]{0.9,0.9,0.9}{
            \makecell[{{p{\linewidth}}}]{
                \texttt{\tiny{[GM$|$GM]}}
                \texttt{COMUNICA: Ciao! Il mio tavolo ha questi oggetti: P(4,4), O(7,6), T(1,3), S(5,6), D(7,7), A(1,3), M(3,2). Nota che A e T sono entrambi in riga 3, colonne 1 e 3. Ho bisogno di allineare tutti gli oggetti con te. Per favore, descrivi la posizione dei tuoi oggetti e proporremo una configurazione obiettivo comune. Potremmo mirare a sistemarli in una riga o colonna libera, ad esempio la riga 5 o 6 che ha più spazio vuoto. Aspetto la tua risposta!} \\
            }
        }
    }
    & & \\ \\

    \theutterance \stepcounter{utterance}  
    & & \multicolumn{4}{p{0.6\linewidth}}{
        \cellcolor[rgb]{0.9,0.9,0.9}{
            \makecell[{{p{\linewidth}}}]{
                \texttt{\tiny{[GM$\rangle$P2]}}
                \texttt{Sono il tuo game master, e stai giocando un gioco collaborativo con la seguente griglia come tavolo da gioco:} \\
\\ 
\texttt{```} \\
\texttt{ 1234567} \\
\texttt{╔═══╤═╤═╗} \\
\texttt{║◌MS│◌│◌║ 1} \\
\texttt{║◌O◌│◌├─╢ 2} \\
\texttt{║◌◌◌├─┘◌║ 3} \\
\texttt{║◌◌◌│◌◌◌║ 4} \\
\texttt{╟─┬─┼───╢ 5} \\
\texttt{║◌│◌│◌◌D║ 6} \\
\texttt{║T│◌│◌PA║ 7} \\
\texttt{╚═╧═╧═══╝} \\
\\ 
\texttt{```} \\
\\ 
\texttt{* Il bordo superiore mostra le coordinate x che aumentano verso destra, e il bordo destro le coordinate y che aumentano verso il basso.} \\
\texttt{* The following objects are randomly placed on your grid: 'P', 'O', 'T', 'S', 'D', 'A', 'M'.} \\
\\ 
\\ 
\texttt{L'altro giocatore vede una variante del tavolo da gioco, dove gli oggetti sono posizionati in diverse posizioni casuali. Non puoi vedere il tavolo dell'altro giocatore, e loro non possono vedere il tuo.} \\
\\ 
\\ 
\texttt{**Obiettivo:**} \\
\\ 
\texttt{Entrambi i giocatori devono spostare gli oggetti sui rispettivi sfondi in modo che oggetti identici finiscano alle stesse coordinate. Devi comunicare con l'altro giocatore per concordare una configurazione obiettivo comune.} \\
\\ 
\\ 
\texttt{**Regole:**} \\
\\ 
\texttt{* In ogni turno, puoi inviare esattamente uno dei seguenti due comandi:} \\
\texttt{ 1. `COMUNICA: <MESSAGE>`: per inviare un messaggio (tutto fino alla prossima interruzione di riga) all'altro giocatore. Lo inoltrerò al tuo partner.} \\
\texttt{2. `SPOSTA: <OBJECT>, (<X>, <Y>)`: per spostare un oggetto in una nuova posizione, dove <X> è la colonna e <Y> è la riga. Ti farò sapere se la tua mossa è andata a buon fine o no.} \\
\texttt{* Se non rispetti il formato, o invii diversi comandi contemporaneamente, devo penalizzarti.} \\
\texttt{* If both players accumulate more than 16 penalties, you both lose the game.} \\
\texttt{* È vitale che comunichi con l'altro giocatore riguardo al tuo stato obiettivo! L'*unico* modo in cui puoi trasmettere la tua strategia all'altro giocatore è usando il comando `COMUNICA: <MESSAGE>`!} \\
\\ 
\\ 
\texttt{**Spostamento Oggetti**} \\
\\ 
\texttt{* Puoi spostare oggetti solo in celle entro i limiti della griglia. La cella di destinazione deve essere vuota, cioè deve contenere solo il simbolo '◌'.} \\
\texttt{* Se provi a spostare un oggetto in un posto che non è vuoto, o provi a spostarlo fuori dalla griglia, devo penalizzarti. Hai un altro tentativo.} \\
\texttt{* Prima di fare una mossa, controlla due volte che il posto di destinazione sia vuoto, e non contenga alcuna lettera, cornice, o linea!} \\
\\ 
\\ 
\texttt{**Fine del Gioco**} \\
\\ 
\texttt{Se pensi di aver raggiunto l'obiettivo di allineare tutti gli oggetti, puoi chiedere all'altro giocatore di finire il gioco inviando COMUNICA: finito?. Se l'altro giocatore ti chiede di finire il gioco, e rispondi con COMUNICA: finito!, the game will end.} \\
\\ 
\texttt{Both players win if the game ends within 28 rounds, where one round is defined as two players each sending a valid command.} \\
\\ 
\texttt{**Punteggio:**} \\
\\ 
\texttt{Più vicini sono gli oggetti identici in entrambi i tavoli da gioco, più punti ottieni. Le penalità riducono i tuoi punti. Riesci a battere il record?} \\
\\ 
\texttt{L'altro giocatore ha iniziato il gioco inviando questo messaggio:} \\
\\ 
\texttt{"Ciao! Il mio tavolo ha questi oggetti: P(4,4), O(7,6), T(1,3), S(5,6), D(7,7), A(1,3), M(3,2). Nota che A e T sono entrambi in riga 3, colonne 1 e 3. Ho bisogno di allineare tutti gli oggetti con te. Per favore, descrivi la posizione dei tuoi oggetti e proporremo una configurazione obiettivo comune. Potremmo mirare a sistemarli in una riga o colonna libera, ad esempio la riga 5 o 6 che ha più spazio vuoto. Aspetto la tua risposta!"} \\
\\ 
\texttt{Qual è il tuo primo comando?} \\
            }
        }
    }
    & \\ \\

    \theutterance \stepcounter{utterance}  
    & & & \multicolumn{4}{p{0.6\linewidth}}{
        \cellcolor[rgb]{0.9,0.9,0.9}{
            \makecell[{{p{\linewidth}}}]{
                \texttt{\tiny{[GM$\langle$P2]}}
                \texttt{COMUNICA: Ciao! Il mio tavolo ha questi oggetti: P(7,7), O(2,2), T(1,7), S(2,1), D(7,6), A(7,7), M(1,1). Nota che P e A sono entrambi in (7,7) — probabilmente un errore di lettura, ma sulla mia griglia A è in (7,7) e P è anch'esso in (7,7), il che è impossibile. Aspetta: ricontrollando la griglia...  } \\
\\ 
\texttt{Eccoti la mia griglia riletta con precisione (coordinate X,Y):  } \\
\texttt{{-} M è in (1,1) — prima riga, prima colonna (dove c'è "MS", quindi M in (1,1), S in (2,1))  } \\
\texttt{{-} S è in (2,1)  } \\
\texttt{{-} O è in (2,2) — seconda riga, seconda colonna (dove c'è "O◌")  } \\
\texttt{{-} T è in (1,7) — settima riga, prima colonna (T in (1,7))  } \\
\texttt{{-} D è in (7,6) — sesta riga, settima colonna  } \\
\texttt{{-} A è in (7,7) — settima riga, settima colonna (PA → P in} \\
            }
        }
    }
     \\ \\

    \theutterance \stepcounter{utterance}  
    & & & \multicolumn{2}{p{0.3\linewidth}}{
        \cellcolor[rgb]{0.9,0.9,0.9}{
            \makecell[{{p{\linewidth}}}]{
                \texttt{\tiny{[GM$|$GM]}}
                \texttt{COMUNICA: Ciao! Il mio tavolo ha questi oggetti: P(7,7), O(2,2), T(1,7), S(2,1), D(7,6), A(7,7), M(1,1). Nota che P e A sono entrambi in (7,7) — probabilmente un errore di lettura, ma sulla mia griglia A è in (7,7) e P è anch'esso in (7,7), il che è impossibile. Aspetta: ricontrollando la griglia...  } \\
\\ 
\texttt{Eccoti la mia griglia riletta con precisione (coordinate X,Y):  } \\
\texttt{{-} M è in (1,1) — prima riga, prima colonna (dove c'è "MS", quindi M in (1,1), S in (2,1))  } \\
\texttt{{-} S è in (2,1)  } \\
\texttt{{-} O è in (2,2) — seconda riga, seconda colonna (dove c'è "O◌")  } \\
\texttt{{-} T è in (1,7) — settima riga, prima colonna (T in (1,7))  } \\
\texttt{{-} D è in (7,6) — sesta riga, settima colonna  } \\
\texttt{{-} A è in (7,7) — settima riga, settima colonna (PA → P in} \\
            }
        }
    }
    & & \\ \\

    \theutterance \stepcounter{utterance}  
    & & & \multicolumn{2}{p{0.3\linewidth}}{
        \cellcolor[rgb]{0.9,0.9,0.9}{
            \makecell[{{p{\linewidth}}}]{
                \texttt{\tiny{[GM$|$GM]}}
                \texttt{Invalid format: tail is not empty: '} \\
\\ 
\texttt{Eccoti la mia griglia riletta con precisione (coordinate X,Y):  } \\
\texttt{{-} M è in (1,1) — prima riga, prima colonna (dove c'è "MS", quindi M in (1,1), S in (2,1))  } \\
\texttt{{-} S è in (2,1)  } \\
\texttt{{-} O è in (2,2) — seconda riga, seconda colonna (dove c'è "O◌")  } \\
\texttt{{-} T è in (1,7) — settima riga, prima colonna (T in (1,7))  } \\
\texttt{{-} D è in (7,6) — sesta riga, settima colonna  } \\
\texttt{{-} A è in (7,7) — settima riga, settima colonna (PA → P in} \\
            }
        }
    }
    & & \\ \\

    \theutterance \stepcounter{utterance}  
    & & \multicolumn{4}{p{0.6\linewidth}}{
        \cellcolor[rgb]{0.9,0.9,0.9}{
            \makecell[{{p{\linewidth}}}]{
                \texttt{\tiny{[GM$\rangle$P2]}}
                \texttt{Penalità: Il tuo messaggio non deve contenere nulla dopo il comando!} \\
\\ 
\texttt{Assicurati che la tua risposta contenga solo `COMUNICA: <MESSAGE>` o `SPOSTA: <OBJECT>, (<X>, <Y>)`, e nient'altro!} \\
\\ 
\texttt{You have collectively accumulated 1 of 16 penalties.} \\
\texttt{Riprova!} \\
            }
        }
    }
    & \\ \\

    \theutterance \stepcounter{utterance}  
    & & & \multicolumn{4}{p{0.6\linewidth}}{
        \cellcolor[rgb]{0.9,0.9,0.9}{
            \makecell[{{p{\linewidth}}}]{
                \texttt{\tiny{[GM$\langle$P2]}}
                \texttt{COMUNICA: Ciao! Il mio tavolo ha questi oggetti: M(1,1), S(2,1), O(2,2), T(1,7), D(7,6), A(7,7), P(7,7). Nota che P e A sono entrambi in (7,7) — questo è un conflitto sulla mia griglia, ma guardando la riga 7: "T│◌│◌PA" → P è in (6,7) e A in (7,7). Correggo: P(6,7), A(7,7). Tutti gli altri sono corretti. Propongo di muovere tutti gli oggetti verso la riga 5 che è quasi vuota e ha spazio per allinearli in colonna. Che ne pensi?} \\
            }
        }
    }
     \\ \\

    \theutterance \stepcounter{utterance}  
    & & & \multicolumn{2}{p{0.3\linewidth}}{
        \cellcolor[rgb]{0.9,0.9,0.9}{
            \makecell[{{p{\linewidth}}}]{
                \texttt{\tiny{[GM$|$GM]}}
                \texttt{COMUNICA: Ciao! Il mio tavolo ha questi oggetti: M(1,1), S(2,1), O(2,2), T(1,7), D(7,6), A(7,7), P(7,7). Nota che P e A sono entrambi in (7,7) — questo è un conflitto sulla mia griglia, ma guardando la riga 7: "T│◌│◌PA" → P è in (6,7) e A in (7,7). Correggo: P(6,7), A(7,7). Tutti gli altri sono corretti. Propongo di muovere tutti gli oggetti verso la riga 5 che è quasi vuota e ha spazio per allinearli in colonna. Che ne pensi?} \\
            }
        }
    }
    & & \\ \\

    \theutterance \stepcounter{utterance}  
    & & \multicolumn{4}{p{0.6\linewidth}}{
        \cellcolor[rgb]{0.9,0.9,0.9}{
            \makecell[{{p{\linewidth}}}]{
                \texttt{\tiny{[P1$\langle$GM]}}
                \texttt{Il tuo messaggio è stato inoltrato all'altro giocatore.} \\
\\ 
\\ 
\texttt{You are currently playing round 1 of maximum 28.} \\
\texttt{You have collectively accumulated 1 of 16 penalties.} \\
\texttt{L'altro giocatore ha inviato questo messaggio:} \\
\\ 
\\ 
\texttt{"Ciao! Il mio tavolo ha questi oggetti: M(1,1), S(2,1), O(2,2), T(1,7), D(7,6), A(7,7), P(7,7). Nota che P e A sono entrambi in (7,7) — questo è un conflitto sulla mia griglia, ma guardando la riga 7: "T│◌│◌PA" → P è in (6,7) e A in (7,7). Correggo: P(6,7), A(7,7). Tutti gli altri sono corretti. Propongo di muovere tutti gli oggetti verso la riga 5 che è quasi vuota e ha spazio per allinearli in colonna. Che ne pensi?"} \\
\\ 
\\ 
\texttt{Qual è il tuo prossimo comando?} \\
            }
        }
    }
    & \\ \\

    \theutterance \stepcounter{utterance}  
    & \multicolumn{4}{p{0.6\linewidth}}{
        \cellcolor[rgb]{0.9,0.9,0.9}{
            \makecell[{{p{\linewidth}}}]{
                \texttt{\tiny{[P1$\rangle$GM]}}
                \texttt{COMUNICA: Grazie per la correzione! Confermo la tua mappa: M(1,1), S(2,1), O(2,2), T(1,7), D(7,6), P(6,7), A(7,7).  } \\
\texttt{La proposta della riga 5 è ottima — è la più libera e permette allineamenti verticali/orizontali. Propongo questa configurazione obiettivo comune:  } \\
\\ 
\texttt{**Obiettivo Comune (Riga 5):**  } \\
\texttt{{-} A → (1,5)  } \\
\texttt{{-} M → (2,5)  } \\
\texttt{{-} S → (3,5)  } \\
\texttt{{-} O → (4,5)  } \\
\texttt{{-} T → (5,5)  } \\
\texttt{{-} P → (6,5)  } \\
\texttt{{-} D → (7,5)  } \\
\\ 
\texttt{Tutti in riga 5, da sinistra a destra, in ordine alfabetico. È pulito, fattibile (la riga 5 è: "◌│◌│◌│◌│◌│◌│◌"), e facile da ricordare.  } \\
\texttt{Se sei d'accordo, iniziamo a muovere gli oggetti verso queste posizioni.  } \\
\texttt{Dichiaro: pronto a spostare il primo oggetto non appena mi confermi “confermato obiettivo”.} \\
            }
        }
    }
    & & \\ \\

    \theutterance \stepcounter{utterance}  
    & & & \multicolumn{2}{p{0.3\linewidth}}{
        \cellcolor[rgb]{0.9,0.9,0.9}{
            \makecell[{{p{\linewidth}}}]{
                \texttt{\tiny{[GM$|$GM]}}
                \texttt{COMUNICA: Grazie per la correzione! Confermo la tua mappa: M(1,1), S(2,1), O(2,2), T(1,7), D(7,6), P(6,7), A(7,7).  } \\
\texttt{La proposta della riga 5 è ottima — è la più libera e permette allineamenti verticali/orizontali. Propongo questa configurazione obiettivo comune:  } \\
\\ 
\texttt{**Obiettivo Comune (Riga 5):**  } \\
\texttt{{-} A → (1,5)  } \\
\texttt{{-} M → (2,5)  } \\
\texttt{{-} S → (3,5)  } \\
\texttt{{-} O → (4,5)  } \\
\texttt{{-} T → (5,5)  } \\
\texttt{{-} P → (6,5)  } \\
\texttt{{-} D → (7,5)  } \\
\\ 
\texttt{Tutti in riga 5, da sinistra a destra, in ordine alfabetico. È pulito, fattibile (la riga 5 è: "◌│◌│◌│◌│◌│◌│◌"), e facile da ricordare.  } \\
\texttt{Se sei d'accordo, iniziamo a muovere gli oggetti verso queste posizioni.  } \\
\texttt{Dichiaro: pronto a spostare il primo oggetto non appena mi confermi “confermato obiettivo”.} \\
            }
        }
    }
    & & \\ \\

    \theutterance \stepcounter{utterance}  
    & & & \multicolumn{2}{p{0.3\linewidth}}{
        \cellcolor[rgb]{0.9,0.9,0.9}{
            \makecell[{{p{\linewidth}}}]{
                \texttt{\tiny{[GM$|$GM]}}
                \texttt{Invalid format: tail is not empty: '} \\
\texttt{La proposta della riga 5 è ottima — è la più libera e permette allineamenti verticali/orizontali. Propongo questa configurazione obiettivo comune:  } \\
\\ 
\texttt{**Obiettivo Comune (Riga 5):**  } \\
\texttt{{-} A → (1,5)  } \\
\texttt{{-} M → (2,5)  } \\
\texttt{{-} S → (3,5)  } \\
\texttt{{-} O → (4,5)  } \\
\texttt{{-} T → (5,5)  } \\
\texttt{{-} P → (6,5)  } \\
\texttt{{-} D → (7,5)  } \\
\\ 
\texttt{Tutti in riga 5, da sinistra a destra, in ordine alfabetico. È pulito, fattibile (la riga 5 è: "◌│◌│◌│◌│◌│◌│◌"), e facile da ricordare.  } \\
\texttt{Se sei d'accordo, iniziamo a muovere gli oggetti verso queste posizioni.  } \\
\texttt{Dichiaro: pronto a spostare il primo oggetto non appena mi confermi “confermato obiettivo”.} \\
            }
        }
    }
    & & \\ \\

    \theutterance \stepcounter{utterance}  
    & & \multicolumn{4}{p{0.6\linewidth}}{
        \cellcolor[rgb]{0.9,0.9,0.9}{
            \makecell[{{p{\linewidth}}}]{
                \texttt{\tiny{[P1$\langle$GM]}}
                \texttt{Penalità: Il tuo messaggio non deve contenere nulla dopo il comando!} \\
\\ 
\texttt{Assicurati che la tua risposta contenga solo `COMUNICA: <MESSAGE>` o `SPOSTA: <OBJECT>, (<X>, <Y>)`, e nient'altro!} \\
\\ 
\texttt{You have collectively accumulated 2 of 16 penalties.} \\
\texttt{Riprova!} \\
            }
        }
    }
    & \\ \\

    \theutterance \stepcounter{utterance}  
    & \multicolumn{4}{p{0.6\linewidth}}{
        \cellcolor[rgb]{0.9,0.9,0.9}{
            \makecell[{{p{\linewidth}}}]{
                \texttt{\tiny{[P1$\rangle$GM]}}
                \texttt{COMUNICA: Grazie per la correzione! Confermo la tua mappa: M(1,1), S(2,1), O(2,2), T(1,7), D(7,6), P(6,7), A(7,7). Propongo questa configurazione obiettivo comune: tutti gli oggetti alla riga 5, in ordine alfabetico da sinistra a destra: A(1,5), M(2,5), S(3,5), O(4,5), T(5,5), P(6,5), D(7,5). La riga 5 è completamente vuota e perfetta. Se sei d'accordo, rispondi con "confermato obiettivo".} \\
            }
        }
    }
    & & \\ \\

    \theutterance \stepcounter{utterance}  
    & & & \multicolumn{2}{p{0.3\linewidth}}{
        \cellcolor[rgb]{0.9,0.9,0.9}{
            \makecell[{{p{\linewidth}}}]{
                \texttt{\tiny{[GM$|$GM]}}
                \texttt{COMUNICA: Grazie per la correzione! Confermo la tua mappa: M(1,1), S(2,1), O(2,2), T(1,7), D(7,6), P(6,7), A(7,7). Propongo questa configurazione obiettivo comune: tutti gli oggetti alla riga 5, in ordine alfabetico da sinistra a destra: A(1,5), M(2,5), S(3,5), O(4,5), T(5,5), P(6,5), D(7,5). La riga 5 è completamente vuota e perfetta. Se sei d'accordo, rispondi con "confermato obiettivo".} \\
            }
        }
    }
    & & \\ \\

    \theutterance \stepcounter{utterance}  
    & & \multicolumn{4}{p{0.6\linewidth}}{
        \cellcolor[rgb]{0.9,0.9,0.9}{
            \makecell[{{p{\linewidth}}}]{
                \texttt{\tiny{[GM$\rangle$P2]}}
                \texttt{Il tuo messaggio è stato inoltrato all'altro giocatore.} \\
\\ 
\\ 
\texttt{You are currently playing round 2 of maximum 28.} \\
\texttt{You have collectively accumulated 2 of 16 penalties.} \\
\texttt{L'altro giocatore ha inviato questo messaggio:} \\
\\ 
\\ 
\texttt{"Grazie per la correzione! Confermo la tua mappa: M(1,1), S(2,1), O(2,2), T(1,7), D(7,6), P(6,7), A(7,7). Propongo questa configurazione obiettivo comune: tutti gli oggetti alla riga 5, in ordine alfabetico da sinistra a destra: A(1,5), M(2,5), S(3,5), O(4,5), T(5,5), P(6,5), D(7,5). La riga 5 è completamente vuota e perfetta. Se sei d'accordo, rispondi con "confermato obiettivo"."} \\
\\ 
\\ 
\texttt{Qual è il tuo prossimo comando?} \\
            }
        }
    }
    & \\ \\

    \theutterance \stepcounter{utterance}  
    & & & \multicolumn{4}{p{0.6\linewidth}}{
        \cellcolor[rgb]{0.9,0.9,0.9}{
            \makecell[{{p{\linewidth}}}]{
                \texttt{\tiny{[GM$\langle$P2]}}
                \texttt{COMUNICA: confermato obiettivo} \\
            }
        }
    }
     \\ \\

    \theutterance \stepcounter{utterance}  
    & & & \multicolumn{2}{p{0.3\linewidth}}{
        \cellcolor[rgb]{0.9,0.9,0.9}{
            \makecell[{{p{\linewidth}}}]{
                \texttt{\tiny{[GM$|$GM]}}
                \texttt{COMUNICA: confermato obiettivo} \\
            }
        }
    }
    & & \\ \\

    \theutterance \stepcounter{utterance}  
    & & \multicolumn{4}{p{0.6\linewidth}}{
        \cellcolor[rgb]{0.9,0.9,0.9}{
            \makecell[{{p{\linewidth}}}]{
                \texttt{\tiny{[P1$\langle$GM]}}
                \texttt{Il tuo messaggio è stato inoltrato all'altro giocatore.} \\
\\ 
\\ 
\texttt{You are currently playing round 2 of maximum 28.} \\
\texttt{You have collectively accumulated 2 of 16 penalties.} \\
\texttt{L'altro giocatore ha inviato questo messaggio:} \\
\\ 
\\ 
\texttt{"confermato obiettivo"} \\
\\ 
\\ 
\texttt{Qual è il tuo prossimo comando?} \\
            }
        }
    }
    & \\ \\

    \theutterance \stepcounter{utterance}  
    & \multicolumn{4}{p{0.6\linewidth}}{
        \cellcolor[rgb]{0.9,0.9,0.9}{
            \makecell[{{p{\linewidth}}}]{
                \texttt{\tiny{[P1$\rangle$GM]}}
                \texttt{SPOSTA: A, (1,5)} \\
            }
        }
    }
    & & \\ \\

    \theutterance \stepcounter{utterance}  
    & & & \multicolumn{2}{p{0.3\linewidth}}{
        \cellcolor[rgb]{0.9,0.9,0.9}{
            \makecell[{{p{\linewidth}}}]{
                \texttt{\tiny{[GM$|$GM]}}
                \texttt{SPOSTA: A, (1,5)} \\
            }
        }
    }
    & & \\ \\

    \theutterance \stepcounter{utterance}  
    & & & \multicolumn{2}{p{0.3\linewidth}}{
        \cellcolor[rgb]{0.9,0.9,0.9}{
            \makecell[{{p{\linewidth}}}]{
                \texttt{\tiny{[GM$|$GM]}}
                \texttt{Penalità: (1,5) non è vuoto, ma contiene '─'.} \\
\\ 
\texttt{You have collectively accumulated 3 of 16 penalties.} \\
\texttt{Riprova!} \\
            }
        }
    }
    & & \\ \\

    \theutterance \stepcounter{utterance}  
    & & \multicolumn{4}{p{0.6\linewidth}}{
        \cellcolor[rgb]{0.9,0.9,0.9}{
            \makecell[{{p{\linewidth}}}]{
                \texttt{\tiny{[P1$\langle$GM]}}
                \texttt{Penalità: (1,5) non è vuoto, ma contiene '─'.} \\
\\ 
\texttt{You have collectively accumulated 3 of 16 penalties.} \\
\texttt{Riprova!} \\
            }
        }
    }
    & \\ \\

    \theutterance \stepcounter{utterance}  
    & \multicolumn{4}{p{0.6\linewidth}}{
        \cellcolor[rgb]{0.9,0.9,0.9}{
            \makecell[{{p{\linewidth}}}]{
                \texttt{\tiny{[P1$\rangle$GM]}}
                \texttt{SPOSTA: M, (2,5)} \\
            }
        }
    }
    & & \\ \\

    \theutterance \stepcounter{utterance}  
    & & & \multicolumn{2}{p{0.3\linewidth}}{
        \cellcolor[rgb]{0.9,0.9,0.9}{
            \makecell[{{p{\linewidth}}}]{
                \texttt{\tiny{[GM$|$GM]}}
                \texttt{SPOSTA: M, (2,5)} \\
            }
        }
    }
    & & \\ \\

    \theutterance \stepcounter{utterance}  
    & & & \multicolumn{2}{p{0.3\linewidth}}{
        \cellcolor[rgb]{0.9,0.9,0.9}{
            \makecell[{{p{\linewidth}}}]{
                \texttt{\tiny{[GM$|$GM]}}
                \texttt{Penalità: (2,5) non è vuoto, ma contiene '┬'.} \\
\\ 
\texttt{You have collectively accumulated 4 of 16 penalties.} \\
\texttt{Riprova!} \\
            }
        }
    }
    & & \\ \\

    \theutterance \stepcounter{utterance}  
    & & \multicolumn{4}{p{0.6\linewidth}}{
        \cellcolor[rgb]{0.9,0.9,0.9}{
            \makecell[{{p{\linewidth}}}]{
                \texttt{\tiny{[P1$\langle$GM]}}
                \texttt{Penalità: (2,5) non è vuoto, ma contiene '┬'.} \\
\\ 
\texttt{You have collectively accumulated 4 of 16 penalties.} \\
\texttt{Riprova!} \\
            }
        }
    }
    & \\ \\

    \theutterance \stepcounter{utterance}  
    & \multicolumn{4}{p{0.6\linewidth}}{
        \cellcolor[rgb]{0.9,0.9,0.9}{
            \makecell[{{p{\linewidth}}}]{
                \texttt{\tiny{[P1$\rangle$GM]}}
                \texttt{SPOSTA: P, (6,4)} \\
            }
        }
    }
    & & \\ \\

    \theutterance \stepcounter{utterance}  
    & & & \multicolumn{2}{p{0.3\linewidth}}{
        \cellcolor[rgb]{0.9,0.9,0.9}{
            \makecell[{{p{\linewidth}}}]{
                \texttt{\tiny{[GM$|$GM]}}
                \texttt{SPOSTA: P, (6,4)} \\
            }
        }
    }
    & & \\ \\

    \theutterance \stepcounter{utterance}  
    & & & \multicolumn{2}{p{0.3\linewidth}}{
        \cellcolor[rgb]{0.9,0.9,0.9}{
            \makecell[{{p{\linewidth}}}]{
                \texttt{\tiny{[GM$|$GM]}}
                \texttt{Spostato 'P' in (6,4) con successo. La tua griglia aggiornata appare così:} \\
\texttt{```} \\
\texttt{ 1234567} \\
\texttt{╔═══╤═╤═╗} \\
\texttt{║◌◌◌│◌│◌║ 1} \\
\texttt{║◌◌M│◌├─╢ 2} \\
\texttt{║AT◌├─┘◌║ 3} \\
\texttt{║◌◌◌│◌P◌║ 4} \\
\texttt{╟─┬─┼───╢ 5} \\
\texttt{║◌│◌│S◌O║ 6} \\
\texttt{║◌│◌│◌D◌║ 7} \\
\texttt{╚═╧═╧═══╝} \\
\texttt{```} \\
            }
        }
    }
    & & \\ \\

    \theutterance \stepcounter{utterance}  
    & & \multicolumn{4}{p{0.6\linewidth}}{
        \cellcolor[rgb]{0.9,0.9,0.9}{
            \makecell[{{p{\linewidth}}}]{
                \texttt{\tiny{[GM$\rangle$P2]}}
                \texttt{Il tuo messaggio è stato inoltrato all'altro giocatore.} \\
\\ 
\\ 
\texttt{You are currently playing round 3 of maximum 28.} \\
\texttt{You have collectively accumulated 4 of 16 penalties.} \\
\texttt{L'altro giocatore ha spostato un oggetto sulla sua griglia.} \\
\\ 
\\ 
\texttt{Qual è il tuo prossimo comando?} \\
            }
        }
    }
    & \\ \\

    \theutterance \stepcounter{utterance}  
    & & & \multicolumn{4}{p{0.6\linewidth}}{
        \cellcolor[rgb]{0.9,0.9,0.9}{
            \makecell[{{p{\linewidth}}}]{
                \texttt{\tiny{[GM$\langle$P2]}}
                \texttt{SPOSTA: M, (2,5)} \\
            }
        }
    }
     \\ \\

    \theutterance \stepcounter{utterance}  
    & & & \multicolumn{2}{p{0.3\linewidth}}{
        \cellcolor[rgb]{0.9,0.9,0.9}{
            \makecell[{{p{\linewidth}}}]{
                \texttt{\tiny{[GM$|$GM]}}
                \texttt{SPOSTA: M, (2,5)} \\
            }
        }
    }
    & & \\ \\

    \theutterance \stepcounter{utterance}  
    & & & \multicolumn{2}{p{0.3\linewidth}}{
        \cellcolor[rgb]{0.9,0.9,0.9}{
            \makecell[{{p{\linewidth}}}]{
                \texttt{\tiny{[GM$|$GM]}}
                \texttt{Penalità: (2,5) non è vuoto, ma contiene '┬'.} \\
\\ 
\texttt{You have collectively accumulated 5 of 16 penalties.} \\
\texttt{Riprova!} \\
            }
        }
    }
    & & \\ \\

    \theutterance \stepcounter{utterance}  
    & & \multicolumn{4}{p{0.6\linewidth}}{
        \cellcolor[rgb]{0.9,0.9,0.9}{
            \makecell[{{p{\linewidth}}}]{
                \texttt{\tiny{[GM$\rangle$P2]}}
                \texttt{Penalità: (2,5) non è vuoto, ma contiene '┬'.} \\
\\ 
\texttt{You have collectively accumulated 5 of 16 penalties.} \\
\texttt{Riprova!} \\
            }
        }
    }
    & \\ \\

    \theutterance \stepcounter{utterance}  
    & & & \multicolumn{4}{p{0.6\linewidth}}{
        \cellcolor[rgb]{0.9,0.9,0.9}{
            \makecell[{{p{\linewidth}}}]{
                \texttt{\tiny{[GM$\langle$P2]}}
                \texttt{SPOSTA: M, (1,5)} \\
            }
        }
    }
     \\ \\

    \theutterance \stepcounter{utterance}  
    & & & \multicolumn{2}{p{0.3\linewidth}}{
        \cellcolor[rgb]{0.9,0.9,0.9}{
            \makecell[{{p{\linewidth}}}]{
                \texttt{\tiny{[GM$|$GM]}}
                \texttt{SPOSTA: M, (1,5)} \\
            }
        }
    }
    & & \\ \\

    \theutterance \stepcounter{utterance}  
    & & & \multicolumn{2}{p{0.3\linewidth}}{
        \cellcolor[rgb]{0.9,0.9,0.9}{
            \makecell[{{p{\linewidth}}}]{
                \texttt{\tiny{[GM$|$GM]}}
                \texttt{Penalità: (1,5) non è vuoto, ma contiene '─'.} \\
\\ 
\texttt{You have collectively accumulated 6 of 16 penalties.} \\
\texttt{Riprova!} \\
            }
        }
    }
    & & \\ \\

    \theutterance \stepcounter{utterance}  
    & & \multicolumn{4}{p{0.6\linewidth}}{
        \cellcolor[rgb]{0.9,0.9,0.9}{
            \makecell[{{p{\linewidth}}}]{
                \texttt{\tiny{[GM$\rangle$P2]}}
                \texttt{Penalità: (1,5) non è vuoto, ma contiene '─'.} \\
\\ 
\texttt{You have collectively accumulated 6 of 16 penalties.} \\
\texttt{Riprova!} \\
            }
        }
    }
    & \\ \\

    \theutterance \stepcounter{utterance}  
    & & & \multicolumn{4}{p{0.6\linewidth}}{
        \cellcolor[rgb]{0.9,0.9,0.9}{
            \makecell[{{p{\linewidth}}}]{
                \texttt{\tiny{[GM$\langle$P2]}}
                \texttt{SPOSTA: M, (3,5)} \\
            }
        }
    }
     \\ \\

    \theutterance \stepcounter{utterance}  
    & & & \multicolumn{2}{p{0.3\linewidth}}{
        \cellcolor[rgb]{0.9,0.9,0.9}{
            \makecell[{{p{\linewidth}}}]{
                \texttt{\tiny{[GM$|$GM]}}
                \texttt{SPOSTA: M, (3,5)} \\
            }
        }
    }
    & & \\ \\

    \theutterance \stepcounter{utterance}  
    & & & \multicolumn{2}{p{0.3\linewidth}}{
        \cellcolor[rgb]{0.9,0.9,0.9}{
            \makecell[{{p{\linewidth}}}]{
                \texttt{\tiny{[GM$|$GM]}}
                \texttt{Penalità: (3,5) non è vuoto, ma contiene '─'.} \\
\\ 
\texttt{You have collectively accumulated 7 of 16 penalties.} \\
\texttt{Riprova!} \\
            }
        }
    }
    & & \\ \\

    \theutterance \stepcounter{utterance}  
    & & \multicolumn{4}{p{0.6\linewidth}}{
        \cellcolor[rgb]{0.9,0.9,0.9}{
            \makecell[{{p{\linewidth}}}]{
                \texttt{\tiny{[GM$\rangle$P2]}}
                \texttt{Penalità: (3,5) non è vuoto, ma contiene '─'.} \\
\\ 
\texttt{You have collectively accumulated 7 of 16 penalties.} \\
\texttt{Riprova!} \\
            }
        }
    }
    & \\ \\

    \theutterance \stepcounter{utterance}  
    & & & \multicolumn{4}{p{0.6\linewidth}}{
        \cellcolor[rgb]{0.9,0.9,0.9}{
            \makecell[{{p{\linewidth}}}]{
                \texttt{\tiny{[GM$\langle$P2]}}
                \texttt{SPOSTA: M, (1,6)} \\
            }
        }
    }
     \\ \\

    \theutterance \stepcounter{utterance}  
    & & & \multicolumn{2}{p{0.3\linewidth}}{
        \cellcolor[rgb]{0.9,0.9,0.9}{
            \makecell[{{p{\linewidth}}}]{
                \texttt{\tiny{[GM$|$GM]}}
                \texttt{SPOSTA: M, (1,6)} \\
            }
        }
    }
    & & \\ \\

    \theutterance \stepcounter{utterance}  
    & & & \multicolumn{2}{p{0.3\linewidth}}{
        \cellcolor[rgb]{0.9,0.9,0.9}{
            \makecell[{{p{\linewidth}}}]{
                \texttt{\tiny{[GM$|$GM]}}
                \texttt{Spostato 'M' in (1,6) con successo. La tua griglia aggiornata appare così:} \\
\texttt{```} \\
\texttt{ 1234567} \\
\texttt{╔═══╤═╤═╗} \\
\texttt{║◌◌S│◌│◌║ 1} \\
\texttt{║◌O◌│◌├─╢ 2} \\
\texttt{║◌◌◌├─┘◌║ 3} \\
\texttt{║◌◌◌│◌◌◌║ 4} \\
\texttt{╟─┬─┼───╢ 5} \\
\texttt{║M│◌│◌◌D║ 6} \\
\texttt{║T│◌│◌PA║ 7} \\
\texttt{╚═╧═╧═══╝} \\
\texttt{```} \\
            }
        }
    }
    & & \\ \\

    \theutterance \stepcounter{utterance}  
    & & \multicolumn{4}{p{0.6\linewidth}}{
        \cellcolor[rgb]{0.9,0.9,0.9}{
            \makecell[{{p{\linewidth}}}]{
                \texttt{\tiny{[P1$\langle$GM]}}
                \texttt{Spostato 'P' in (6,4) con successo. La tua griglia aggiornata appare così:} \\
\texttt{```} \\
\texttt{ 1234567} \\
\texttt{╔═══╤═╤═╗} \\
\texttt{║◌◌◌│◌│◌║ 1} \\
\texttt{║◌◌M│◌├─╢ 2} \\
\texttt{║AT◌├─┘◌║ 3} \\
\texttt{║◌◌◌│◌P◌║ 4} \\
\texttt{╟─┬─┼───╢ 5} \\
\texttt{║◌│◌│S◌O║ 6} \\
\texttt{║◌│◌│◌D◌║ 7} \\
\texttt{╚═╧═╧═══╝} \\
\texttt{```} \\
\texttt{You are currently playing round 3 of maximum 28.} \\
\texttt{You have collectively accumulated 7 of 16 penalties.} \\
\texttt{L'altro giocatore ha spostato un oggetto sulla sua griglia.} \\
\\ 
\\ 
\texttt{Qual è il tuo prossimo comando?} \\
            }
        }
    }
    & \\ \\

    \theutterance \stepcounter{utterance}  
    & \multicolumn{4}{p{0.6\linewidth}}{
        \cellcolor[rgb]{0.9,0.9,0.9}{
            \makecell[{{p{\linewidth}}}]{
                \texttt{\tiny{[P1$\rangle$GM]}}
                \texttt{SPOSTA: T, (5,5)} \\
            }
        }
    }
    & & \\ \\

    \theutterance \stepcounter{utterance}  
    & & & \multicolumn{2}{p{0.3\linewidth}}{
        \cellcolor[rgb]{0.9,0.9,0.9}{
            \makecell[{{p{\linewidth}}}]{
                \texttt{\tiny{[GM$|$GM]}}
                \texttt{SPOSTA: T, (5,5)} \\
            }
        }
    }
    & & \\ \\

    \theutterance \stepcounter{utterance}  
    & & & \multicolumn{2}{p{0.3\linewidth}}{
        \cellcolor[rgb]{0.9,0.9,0.9}{
            \makecell[{{p{\linewidth}}}]{
                \texttt{\tiny{[GM$|$GM]}}
                \texttt{Penalità: (5,5) non è vuoto, ma contiene '─'.} \\
\\ 
\texttt{You have collectively accumulated 8 of 16 penalties.} \\
\texttt{Riprova!} \\
            }
        }
    }
    & & \\ \\

    \theutterance \stepcounter{utterance}  
    & & \multicolumn{4}{p{0.6\linewidth}}{
        \cellcolor[rgb]{0.9,0.9,0.9}{
            \makecell[{{p{\linewidth}}}]{
                \texttt{\tiny{[P1$\langle$GM]}}
                \texttt{Penalità: (5,5) non è vuoto, ma contiene '─'.} \\
\\ 
\texttt{You have collectively accumulated 8 of 16 penalties.} \\
\texttt{Riprova!} \\
            }
        }
    }
    & \\ \\

    \theutterance \stepcounter{utterance}  
    & \multicolumn{4}{p{0.6\linewidth}}{
        \cellcolor[rgb]{0.9,0.9,0.9}{
            \makecell[{{p{\linewidth}}}]{
                \texttt{\tiny{[P1$\rangle$GM]}}
                \texttt{SPOSTA: T, (5,4)} \\
            }
        }
    }
    & & \\ \\

    \theutterance \stepcounter{utterance}  
    & & & \multicolumn{2}{p{0.3\linewidth}}{
        \cellcolor[rgb]{0.9,0.9,0.9}{
            \makecell[{{p{\linewidth}}}]{
                \texttt{\tiny{[GM$|$GM]}}
                \texttt{SPOSTA: T, (5,4)} \\
            }
        }
    }
    & & \\ \\

    \theutterance \stepcounter{utterance}  
    & & & \multicolumn{2}{p{0.3\linewidth}}{
        \cellcolor[rgb]{0.9,0.9,0.9}{
            \makecell[{{p{\linewidth}}}]{
                \texttt{\tiny{[GM$|$GM]}}
                \texttt{Spostato 'T' in (5,4) con successo. La tua griglia aggiornata appare così:} \\
\texttt{```} \\
\texttt{ 1234567} \\
\texttt{╔═══╤═╤═╗} \\
\texttt{║◌◌◌│◌│◌║ 1} \\
\texttt{║◌◌M│◌├─╢ 2} \\
\texttt{║A◌◌├─┘◌║ 3} \\
\texttt{║◌◌◌│TP◌║ 4} \\
\texttt{╟─┬─┼───╢ 5} \\
\texttt{║◌│◌│S◌O║ 6} \\
\texttt{║◌│◌│◌D◌║ 7} \\
\texttt{╚═╧═╧═══╝} \\
\texttt{```} \\
            }
        }
    }
    & & \\ \\

    \theutterance \stepcounter{utterance}  
    & & \multicolumn{4}{p{0.6\linewidth}}{
        \cellcolor[rgb]{0.9,0.9,0.9}{
            \makecell[{{p{\linewidth}}}]{
                \texttt{\tiny{[GM$\rangle$P2]}}
                \texttt{Spostato 'M' in (1,6) con successo. La tua griglia aggiornata appare così:} \\
\texttt{```} \\
\texttt{ 1234567} \\
\texttt{╔═══╤═╤═╗} \\
\texttt{║◌◌S│◌│◌║ 1} \\
\texttt{║◌O◌│◌├─╢ 2} \\
\texttt{║◌◌◌├─┘◌║ 3} \\
\texttt{║◌◌◌│◌◌◌║ 4} \\
\texttt{╟─┬─┼───╢ 5} \\
\texttt{║M│◌│◌◌D║ 6} \\
\texttt{║T│◌│◌PA║ 7} \\
\texttt{╚═╧═╧═══╝} \\
\texttt{```} \\
\texttt{You are currently playing round 4 of maximum 28.} \\
\texttt{You have collectively accumulated 8 of 16 penalties.} \\
\texttt{L'altro giocatore ha spostato un oggetto sulla sua griglia.} \\
\\ 
\\ 
\texttt{Qual è il tuo prossimo comando?} \\
            }
        }
    }
    & \\ \\

    \theutterance \stepcounter{utterance}  
    & & & \multicolumn{4}{p{0.6\linewidth}}{
        \cellcolor[rgb]{0.9,0.9,0.9}{
            \makecell[{{p{\linewidth}}}]{
                \texttt{\tiny{[GM$\langle$P2]}}
                \texttt{SPOSTA: S, (3,5)} \\
            }
        }
    }
     \\ \\

    \theutterance \stepcounter{utterance}  
    & & & \multicolumn{2}{p{0.3\linewidth}}{
        \cellcolor[rgb]{0.9,0.9,0.9}{
            \makecell[{{p{\linewidth}}}]{
                \texttt{\tiny{[GM$|$GM]}}
                \texttt{SPOSTA: S, (3,5)} \\
            }
        }
    }
    & & \\ \\

    \theutterance \stepcounter{utterance}  
    & & & \multicolumn{2}{p{0.3\linewidth}}{
        \cellcolor[rgb]{0.9,0.9,0.9}{
            \makecell[{{p{\linewidth}}}]{
                \texttt{\tiny{[GM$|$GM]}}
                \texttt{Penalità: (3,5) non è vuoto, ma contiene '─'.} \\
\\ 
\texttt{You have collectively accumulated 9 of 16 penalties.} \\
\texttt{Riprova!} \\
            }
        }
    }
    & & \\ \\

    \theutterance \stepcounter{utterance}  
    & & \multicolumn{4}{p{0.6\linewidth}}{
        \cellcolor[rgb]{0.9,0.9,0.9}{
            \makecell[{{p{\linewidth}}}]{
                \texttt{\tiny{[GM$\rangle$P2]}}
                \texttt{Penalità: (3,5) non è vuoto, ma contiene '─'.} \\
\\ 
\texttt{You have collectively accumulated 9 of 16 penalties.} \\
\texttt{Riprova!} \\
            }
        }
    }
    & \\ \\

    \theutterance \stepcounter{utterance}  
    & & & \multicolumn{4}{p{0.6\linewidth}}{
        \cellcolor[rgb]{0.9,0.9,0.9}{
            \makecell[{{p{\linewidth}}}]{
                \texttt{\tiny{[GM$\langle$P2]}}
                \texttt{SPOSTA: S, (2,5)} \\
            }
        }
    }
     \\ \\

    \theutterance \stepcounter{utterance}  
    & & & \multicolumn{2}{p{0.3\linewidth}}{
        \cellcolor[rgb]{0.9,0.9,0.9}{
            \makecell[{{p{\linewidth}}}]{
                \texttt{\tiny{[GM$|$GM]}}
                \texttt{SPOSTA: S, (2,5)} \\
            }
        }
    }
    & & \\ \\

    \theutterance \stepcounter{utterance}  
    & & & \multicolumn{2}{p{0.3\linewidth}}{
        \cellcolor[rgb]{0.9,0.9,0.9}{
            \makecell[{{p{\linewidth}}}]{
                \texttt{\tiny{[GM$|$GM]}}
                \texttt{Penalità: (2,5) non è vuoto, ma contiene '┬'.} \\
\\ 
\texttt{You have collectively accumulated 10 of 16 penalties.} \\
\texttt{Riprova!} \\
            }
        }
    }
    & & \\ \\

    \theutterance \stepcounter{utterance}  
    & & \multicolumn{4}{p{0.6\linewidth}}{
        \cellcolor[rgb]{0.9,0.9,0.9}{
            \makecell[{{p{\linewidth}}}]{
                \texttt{\tiny{[GM$\rangle$P2]}}
                \texttt{Penalità: (2,5) non è vuoto, ma contiene '┬'.} \\
\\ 
\texttt{You have collectively accumulated 10 of 16 penalties.} \\
\texttt{Riprova!} \\
            }
        }
    }
    & \\ \\

    \theutterance \stepcounter{utterance}  
    & & & \multicolumn{4}{p{0.6\linewidth}}{
        \cellcolor[rgb]{0.9,0.9,0.9}{
            \makecell[{{p{\linewidth}}}]{
                \texttt{\tiny{[GM$\langle$P2]}}
                \texttt{SPOSTA: S, (4,5)} \\
            }
        }
    }
     \\ \\

    \theutterance \stepcounter{utterance}  
    & & & \multicolumn{2}{p{0.3\linewidth}}{
        \cellcolor[rgb]{0.9,0.9,0.9}{
            \makecell[{{p{\linewidth}}}]{
                \texttt{\tiny{[GM$|$GM]}}
                \texttt{SPOSTA: S, (4,5)} \\
            }
        }
    }
    & & \\ \\

    \theutterance \stepcounter{utterance}  
    & & & \multicolumn{2}{p{0.3\linewidth}}{
        \cellcolor[rgb]{0.9,0.9,0.9}{
            \makecell[{{p{\linewidth}}}]{
                \texttt{\tiny{[GM$|$GM]}}
                \texttt{Penalità: (4,5) non è vuoto, ma contiene '┼'.} \\
\\ 
\texttt{You have collectively accumulated 11 of 16 penalties.} \\
\texttt{Riprova!} \\
            }
        }
    }
    & & \\ \\

    \theutterance \stepcounter{utterance}  
    & & \multicolumn{4}{p{0.6\linewidth}}{
        \cellcolor[rgb]{0.9,0.9,0.9}{
            \makecell[{{p{\linewidth}}}]{
                \texttt{\tiny{[GM$\rangle$P2]}}
                \texttt{Penalità: (4,5) non è vuoto, ma contiene '┼'.} \\
\\ 
\texttt{You have collectively accumulated 11 of 16 penalties.} \\
\texttt{Riprova!} \\
            }
        }
    }
    & \\ \\

    \theutterance \stepcounter{utterance}  
    & & & \multicolumn{4}{p{0.6\linewidth}}{
        \cellcolor[rgb]{0.9,0.9,0.9}{
            \makecell[{{p{\linewidth}}}]{
                \texttt{\tiny{[GM$\langle$P2]}}
                \texttt{SPOSTA: S, (3,4)} \\
            }
        }
    }
     \\ \\

    \theutterance \stepcounter{utterance}  
    & & & \multicolumn{2}{p{0.3\linewidth}}{
        \cellcolor[rgb]{0.9,0.9,0.9}{
            \makecell[{{p{\linewidth}}}]{
                \texttt{\tiny{[GM$|$GM]}}
                \texttt{SPOSTA: S, (3,4)} \\
            }
        }
    }
    & & \\ \\

    \theutterance \stepcounter{utterance}  
    & & & \multicolumn{2}{p{0.3\linewidth}}{
        \cellcolor[rgb]{0.9,0.9,0.9}{
            \makecell[{{p{\linewidth}}}]{
                \texttt{\tiny{[GM$|$GM]}}
                \texttt{Spostato 'S' in (3,4) con successo. La tua griglia aggiornata appare così:} \\
\texttt{```} \\
\texttt{ 1234567} \\
\texttt{╔═══╤═╤═╗} \\
\texttt{║◌◌◌│◌│◌║ 1} \\
\texttt{║◌O◌│◌├─╢ 2} \\
\texttt{║◌◌◌├─┘◌║ 3} \\
\texttt{║◌◌S│◌◌◌║ 4} \\
\texttt{╟─┬─┼───╢ 5} \\
\texttt{║M│◌│◌◌D║ 6} \\
\texttt{║T│◌│◌PA║ 7} \\
\texttt{╚═╧═╧═══╝} \\
\texttt{```} \\
            }
        }
    }
    & & \\ \\

    \theutterance \stepcounter{utterance}  
    & & \multicolumn{4}{p{0.6\linewidth}}{
        \cellcolor[rgb]{0.9,0.9,0.9}{
            \makecell[{{p{\linewidth}}}]{
                \texttt{\tiny{[P1$\langle$GM]}}
                \texttt{Spostato 'T' in (5,4) con successo. La tua griglia aggiornata appare così:} \\
\texttt{```} \\
\texttt{ 1234567} \\
\texttt{╔═══╤═╤═╗} \\
\texttt{║◌◌◌│◌│◌║ 1} \\
\texttt{║◌◌M│◌├─╢ 2} \\
\texttt{║A◌◌├─┘◌║ 3} \\
\texttt{║◌◌◌│TP◌║ 4} \\
\texttt{╟─┬─┼───╢ 5} \\
\texttt{║◌│◌│S◌O║ 6} \\
\texttt{║◌│◌│◌D◌║ 7} \\
\texttt{╚═╧═╧═══╝} \\
\texttt{```} \\
\texttt{You are currently playing round 4 of maximum 28.} \\
\texttt{You have collectively accumulated 11 of 16 penalties.} \\
\texttt{L'altro giocatore ha spostato un oggetto sulla sua griglia.} \\
\\ 
\\ 
\texttt{Qual è il tuo prossimo comando?} \\
            }
        }
    }
    & \\ \\

    \theutterance \stepcounter{utterance}  
    & \multicolumn{4}{p{0.6\linewidth}}{
        \cellcolor[rgb]{0.9,0.9,0.9}{
            \makecell[{{p{\linewidth}}}]{
                \texttt{\tiny{[P1$\rangle$GM]}}
                \texttt{SPOSTA: M, (3,4)} \\
            }
        }
    }
    & & \\ \\

    \theutterance \stepcounter{utterance}  
    & & & \multicolumn{2}{p{0.3\linewidth}}{
        \cellcolor[rgb]{0.9,0.9,0.9}{
            \makecell[{{p{\linewidth}}}]{
                \texttt{\tiny{[GM$|$GM]}}
                \texttt{SPOSTA: M, (3,4)} \\
            }
        }
    }
    & & \\ \\

    \theutterance \stepcounter{utterance}  
    & & & \multicolumn{2}{p{0.3\linewidth}}{
        \cellcolor[rgb]{0.9,0.9,0.9}{
            \makecell[{{p{\linewidth}}}]{
                \texttt{\tiny{[GM$|$GM]}}
                \texttt{Spostato 'M' in (3,4) con successo. La tua griglia aggiornata appare così:} \\
\texttt{```} \\
\texttt{ 1234567} \\
\texttt{╔═══╤═╤═╗} \\
\texttt{║◌◌◌│◌│◌║ 1} \\
\texttt{║◌◌◌│◌├─╢ 2} \\
\texttt{║A◌◌├─┘◌║ 3} \\
\texttt{║◌◌M│TP◌║ 4} \\
\texttt{╟─┬─┼───╢ 5} \\
\texttt{║◌│◌│S◌O║ 6} \\
\texttt{║◌│◌│◌D◌║ 7} \\
\texttt{╚═╧═╧═══╝} \\
\texttt{```} \\
            }
        }
    }
    & & \\ \\

    \theutterance \stepcounter{utterance}  
    & & \multicolumn{4}{p{0.6\linewidth}}{
        \cellcolor[rgb]{0.9,0.9,0.9}{
            \makecell[{{p{\linewidth}}}]{
                \texttt{\tiny{[GM$\rangle$P2]}}
                \texttt{Spostato 'S' in (3,4) con successo. La tua griglia aggiornata appare così:} \\
\texttt{```} \\
\texttt{ 1234567} \\
\texttt{╔═══╤═╤═╗} \\
\texttt{║◌◌◌│◌│◌║ 1} \\
\texttt{║◌O◌│◌├─╢ 2} \\
\texttt{║◌◌◌├─┘◌║ 3} \\
\texttt{║◌◌S│◌◌◌║ 4} \\
\texttt{╟─┬─┼───╢ 5} \\
\texttt{║M│◌│◌◌D║ 6} \\
\texttt{║T│◌│◌PA║ 7} \\
\texttt{╚═╧═╧═══╝} \\
\texttt{```} \\
\texttt{You are currently playing round 5 of maximum 28.} \\
\texttt{You have collectively accumulated 11 of 16 penalties.} \\
\texttt{L'altro giocatore ha spostato un oggetto sulla sua griglia.} \\
\\ 
\\ 
\texttt{Qual è il tuo prossimo comando?} \\
            }
        }
    }
    & \\ \\

    \theutterance \stepcounter{utterance}  
    & & & \multicolumn{4}{p{0.6\linewidth}}{
        \cellcolor[rgb]{0.9,0.9,0.9}{
            \makecell[{{p{\linewidth}}}]{
                \texttt{\tiny{[GM$\langle$P2]}}
                \texttt{SPOSTA: O, (4,5)} \\
            }
        }
    }
     \\ \\

    \theutterance \stepcounter{utterance}  
    & & & \multicolumn{2}{p{0.3\linewidth}}{
        \cellcolor[rgb]{0.9,0.9,0.9}{
            \makecell[{{p{\linewidth}}}]{
                \texttt{\tiny{[GM$|$GM]}}
                \texttt{SPOSTA: O, (4,5)} \\
            }
        }
    }
    & & \\ \\

    \theutterance \stepcounter{utterance}  
    & & & \multicolumn{2}{p{0.3\linewidth}}{
        \cellcolor[rgb]{0.9,0.9,0.9}{
            \makecell[{{p{\linewidth}}}]{
                \texttt{\tiny{[GM$|$GM]}}
                \texttt{Penalità: (4,5) non è vuoto, ma contiene '┼'.} \\
\\ 
\texttt{You have collectively accumulated 12 of 16 penalties.} \\
\texttt{Riprova!} \\
            }
        }
    }
    & & \\ \\

    \theutterance \stepcounter{utterance}  
    & & \multicolumn{4}{p{0.6\linewidth}}{
        \cellcolor[rgb]{0.9,0.9,0.9}{
            \makecell[{{p{\linewidth}}}]{
                \texttt{\tiny{[GM$\rangle$P2]}}
                \texttt{Penalità: (4,5) non è vuoto, ma contiene '┼'.} \\
\\ 
\texttt{You have collectively accumulated 12 of 16 penalties.} \\
\texttt{Riprova!} \\
            }
        }
    }
    & \\ \\

    \theutterance \stepcounter{utterance}  
    & & & \multicolumn{4}{p{0.6\linewidth}}{
        \cellcolor[rgb]{0.9,0.9,0.9}{
            \makecell[{{p{\linewidth}}}]{
                \texttt{\tiny{[GM$\langle$P2]}}
                \texttt{SPOSTA: O, (4,4)} \\
            }
        }
    }
     \\ \\

    \theutterance \stepcounter{utterance}  
    & & & \multicolumn{2}{p{0.3\linewidth}}{
        \cellcolor[rgb]{0.9,0.9,0.9}{
            \makecell[{{p{\linewidth}}}]{
                \texttt{\tiny{[GM$|$GM]}}
                \texttt{SPOSTA: O, (4,4)} \\
            }
        }
    }
    & & \\ \\

    \theutterance \stepcounter{utterance}  
    & & & \multicolumn{2}{p{0.3\linewidth}}{
        \cellcolor[rgb]{0.9,0.9,0.9}{
            \makecell[{{p{\linewidth}}}]{
                \texttt{\tiny{[GM$|$GM]}}
                \texttt{Penalità: (4,4) non è vuoto, ma contiene '│'.} \\
\\ 
\texttt{You have collectively accumulated 13 of 16 penalties.} \\
\texttt{Riprova!} \\
            }
        }
    }
    & & \\ \\

    \theutterance \stepcounter{utterance}  
    & & \multicolumn{4}{p{0.6\linewidth}}{
        \cellcolor[rgb]{0.9,0.9,0.9}{
            \makecell[{{p{\linewidth}}}]{
                \texttt{\tiny{[GM$\rangle$P2]}}
                \texttt{Penalità: (4,4) non è vuoto, ma contiene '│'.} \\
\\ 
\texttt{You have collectively accumulated 13 of 16 penalties.} \\
\texttt{Riprova!} \\
            }
        }
    }
    & \\ \\

    \theutterance \stepcounter{utterance}  
    & & & \multicolumn{4}{p{0.6\linewidth}}{
        \cellcolor[rgb]{0.9,0.9,0.9}{
            \makecell[{{p{\linewidth}}}]{
                \texttt{\tiny{[GM$\langle$P2]}}
                \texttt{SPOSTA: O, (5,4)} \\
            }
        }
    }
     \\ \\

    \theutterance \stepcounter{utterance}  
    & & & \multicolumn{2}{p{0.3\linewidth}}{
        \cellcolor[rgb]{0.9,0.9,0.9}{
            \makecell[{{p{\linewidth}}}]{
                \texttt{\tiny{[GM$|$GM]}}
                \texttt{SPOSTA: O, (5,4)} \\
            }
        }
    }
    & & \\ \\

    \theutterance \stepcounter{utterance}  
    & & & \multicolumn{2}{p{0.3\linewidth}}{
        \cellcolor[rgb]{0.9,0.9,0.9}{
            \makecell[{{p{\linewidth}}}]{
                \texttt{\tiny{[GM$|$GM]}}
                \texttt{Spostato 'O' in (5,4) con successo. La tua griglia aggiornata appare così:} \\
\texttt{```} \\
\texttt{ 1234567} \\
\texttt{╔═══╤═╤═╗} \\
\texttt{║◌◌◌│◌│◌║ 1} \\
\texttt{║◌◌◌│◌├─╢ 2} \\
\texttt{║◌◌◌├─┘◌║ 3} \\
\texttt{║◌◌S│O◌◌║ 4} \\
\texttt{╟─┬─┼───╢ 5} \\
\texttt{║M│◌│◌◌D║ 6} \\
\texttt{║T│◌│◌PA║ 7} \\
\texttt{╚═╧═╧═══╝} \\
\texttt{```} \\
            }
        }
    }
    & & \\ \\

    \theutterance \stepcounter{utterance}  
    & & \multicolumn{4}{p{0.6\linewidth}}{
        \cellcolor[rgb]{0.9,0.9,0.9}{
            \makecell[{{p{\linewidth}}}]{
                \texttt{\tiny{[P1$\langle$GM]}}
                \texttt{Spostato 'M' in (3,4) con successo. La tua griglia aggiornata appare così:} \\
\texttt{```} \\
\texttt{ 1234567} \\
\texttt{╔═══╤═╤═╗} \\
\texttt{║◌◌◌│◌│◌║ 1} \\
\texttt{║◌◌◌│◌├─╢ 2} \\
\texttt{║A◌◌├─┘◌║ 3} \\
\texttt{║◌◌M│TP◌║ 4} \\
\texttt{╟─┬─┼───╢ 5} \\
\texttt{║◌│◌│S◌O║ 6} \\
\texttt{║◌│◌│◌D◌║ 7} \\
\texttt{╚═╧═╧═══╝} \\
\texttt{```} \\
\texttt{You are currently playing round 5 of maximum 28.} \\
\texttt{You have collectively accumulated 13 of 16 penalties.} \\
\texttt{L'altro giocatore ha spostato un oggetto sulla sua griglia.} \\
\\ 
\\ 
\texttt{Qual è il tuo prossimo comando?} \\
            }
        }
    }
    & \\ \\

    \theutterance \stepcounter{utterance}  
    & \multicolumn{4}{p{0.6\linewidth}}{
        \cellcolor[rgb]{0.9,0.9,0.9}{
            \makecell[{{p{\linewidth}}}]{
                \texttt{\tiny{[P1$\rangle$GM]}}
                \texttt{SPOSTA: A, (1,4)} \\
            }
        }
    }
    & & \\ \\

    \theutterance \stepcounter{utterance}  
    & & & \multicolumn{2}{p{0.3\linewidth}}{
        \cellcolor[rgb]{0.9,0.9,0.9}{
            \makecell[{{p{\linewidth}}}]{
                \texttt{\tiny{[GM$|$GM]}}
                \texttt{SPOSTA: A, (1,4)} \\
            }
        }
    }
    & & \\ \\

    \theutterance \stepcounter{utterance}  
    & & & \multicolumn{2}{p{0.3\linewidth}}{
        \cellcolor[rgb]{0.9,0.9,0.9}{
            \makecell[{{p{\linewidth}}}]{
                \texttt{\tiny{[GM$|$GM]}}
                \texttt{Spostato 'A' in (1,4) con successo. La tua griglia aggiornata appare così:} \\
\texttt{```} \\
\texttt{ 1234567} \\
\texttt{╔═══╤═╤═╗} \\
\texttt{║◌◌◌│◌│◌║ 1} \\
\texttt{║◌◌◌│◌├─╢ 2} \\
\texttt{║◌◌◌├─┘◌║ 3} \\
\texttt{║A◌M│TP◌║ 4} \\
\texttt{╟─┬─┼───╢ 5} \\
\texttt{║◌│◌│S◌O║ 6} \\
\texttt{║◌│◌│◌D◌║ 7} \\
\texttt{╚═╧═╧═══╝} \\
\texttt{```} \\
            }
        }
    }
    & & \\ \\

    \theutterance \stepcounter{utterance}  
    & & \multicolumn{4}{p{0.6\linewidth}}{
        \cellcolor[rgb]{0.9,0.9,0.9}{
            \makecell[{{p{\linewidth}}}]{
                \texttt{\tiny{[GM$\rangle$P2]}}
                \texttt{Spostato 'O' in (5,4) con successo. La tua griglia aggiornata appare così:} \\
\texttt{```} \\
\texttt{ 1234567} \\
\texttt{╔═══╤═╤═╗} \\
\texttt{║◌◌◌│◌│◌║ 1} \\
\texttt{║◌◌◌│◌├─╢ 2} \\
\texttt{║◌◌◌├─┘◌║ 3} \\
\texttt{║◌◌S│O◌◌║ 4} \\
\texttt{╟─┬─┼───╢ 5} \\
\texttt{║M│◌│◌◌D║ 6} \\
\texttt{║T│◌│◌PA║ 7} \\
\texttt{╚═╧═╧═══╝} \\
\texttt{```} \\
\texttt{You are currently playing round 6 of maximum 28.} \\
\texttt{You have collectively accumulated 13 of 16 penalties.} \\
\texttt{L'altro giocatore ha spostato un oggetto sulla sua griglia.} \\
\\ 
\\ 
\texttt{Qual è il tuo prossimo comando?} \\
            }
        }
    }
    & \\ \\

    \theutterance \stepcounter{utterance}  
    & & & \multicolumn{4}{p{0.6\linewidth}}{
        \cellcolor[rgb]{0.9,0.9,0.9}{
            \makecell[{{p{\linewidth}}}]{
                \texttt{\tiny{[GM$\langle$P2]}}
                \texttt{SPOSTA: T, (5,5)} \\
            }
        }
    }
     \\ \\

    \theutterance \stepcounter{utterance}  
    & & & \multicolumn{2}{p{0.3\linewidth}}{
        \cellcolor[rgb]{0.9,0.9,0.9}{
            \makecell[{{p{\linewidth}}}]{
                \texttt{\tiny{[GM$|$GM]}}
                \texttt{SPOSTA: T, (5,5)} \\
            }
        }
    }
    & & \\ \\

    \theutterance \stepcounter{utterance}  
    & & & \multicolumn{2}{p{0.3\linewidth}}{
        \cellcolor[rgb]{0.9,0.9,0.9}{
            \makecell[{{p{\linewidth}}}]{
                \texttt{\tiny{[GM$|$GM]}}
                \texttt{Penalità: (5,5) non è vuoto, ma contiene '─'.} \\
\\ 
\texttt{You have collectively accumulated 14 of 16 penalties.} \\
\texttt{Riprova!} \\
            }
        }
    }
    & & \\ \\

    \theutterance \stepcounter{utterance}  
    & & \multicolumn{4}{p{0.6\linewidth}}{
        \cellcolor[rgb]{0.9,0.9,0.9}{
            \makecell[{{p{\linewidth}}}]{
                \texttt{\tiny{[GM$\rangle$P2]}}
                \texttt{Penalità: (5,5) non è vuoto, ma contiene '─'.} \\
\\ 
\texttt{You have collectively accumulated 14 of 16 penalties.} \\
\texttt{Riprova!} \\
            }
        }
    }
    & \\ \\

    \theutterance \stepcounter{utterance}  
    & & & \multicolumn{4}{p{0.6\linewidth}}{
        \cellcolor[rgb]{0.9,0.9,0.9}{
            \makecell[{{p{\linewidth}}}]{
                \texttt{\tiny{[GM$\langle$P2]}}
                \texttt{SPOSTA: T, (5,6)} \\
            }
        }
    }
     \\ \\

    \theutterance \stepcounter{utterance}  
    & & & \multicolumn{2}{p{0.3\linewidth}}{
        \cellcolor[rgb]{0.9,0.9,0.9}{
            \makecell[{{p{\linewidth}}}]{
                \texttt{\tiny{[GM$|$GM]}}
                \texttt{SPOSTA: T, (5,6)} \\
            }
        }
    }
    & & \\ \\

    \theutterance \stepcounter{utterance}  
    & & & \multicolumn{2}{p{0.3\linewidth}}{
        \cellcolor[rgb]{0.9,0.9,0.9}{
            \makecell[{{p{\linewidth}}}]{
                \texttt{\tiny{[GM$|$GM]}}
                \texttt{Spostato 'T' in (5,6) con successo. La tua griglia aggiornata appare così:} \\
\texttt{```} \\
\texttt{ 1234567} \\
\texttt{╔═══╤═╤═╗} \\
\texttt{║◌◌◌│◌│◌║ 1} \\
\texttt{║◌◌◌│◌├─╢ 2} \\
\texttt{║◌◌◌├─┘◌║ 3} \\
\texttt{║◌◌S│O◌◌║ 4} \\
\texttt{╟─┬─┼───╢ 5} \\
\texttt{║M│◌│T◌D║ 6} \\
\texttt{║◌│◌│◌PA║ 7} \\
\texttt{╚═╧═╧═══╝} \\
\texttt{```} \\
            }
        }
    }
    & & \\ \\

    \theutterance \stepcounter{utterance}  
    & & \multicolumn{4}{p{0.6\linewidth}}{
        \cellcolor[rgb]{0.9,0.9,0.9}{
            \makecell[{{p{\linewidth}}}]{
                \texttt{\tiny{[P1$\langle$GM]}}
                \texttt{Spostato 'A' in (1,4) con successo. La tua griglia aggiornata appare così:} \\
\texttt{```} \\
\texttt{ 1234567} \\
\texttt{╔═══╤═╤═╗} \\
\texttt{║◌◌◌│◌│◌║ 1} \\
\texttt{║◌◌◌│◌├─╢ 2} \\
\texttt{║◌◌◌├─┘◌║ 3} \\
\texttt{║A◌M│TP◌║ 4} \\
\texttt{╟─┬─┼───╢ 5} \\
\texttt{║◌│◌│S◌O║ 6} \\
\texttt{║◌│◌│◌D◌║ 7} \\
\texttt{╚═╧═╧═══╝} \\
\texttt{```} \\
\texttt{You are currently playing round 6 of maximum 28.} \\
\texttt{You have collectively accumulated 14 of 16 penalties.} \\
\texttt{L'altro giocatore ha spostato un oggetto sulla sua griglia.} \\
\\ 
\\ 
\texttt{Qual è il tuo prossimo comando?} \\
            }
        }
    }
    & \\ \\

    \theutterance \stepcounter{utterance}  
    & \multicolumn{4}{p{0.6\linewidth}}{
        \cellcolor[rgb]{0.9,0.9,0.9}{
            \makecell[{{p{\linewidth}}}]{
                \texttt{\tiny{[P1$\rangle$GM]}}
                \texttt{SPOSTA: S, (2,5)} \\
            }
        }
    }
    & & \\ \\

    \theutterance \stepcounter{utterance}  
    & & & \multicolumn{2}{p{0.3\linewidth}}{
        \cellcolor[rgb]{0.9,0.9,0.9}{
            \makecell[{{p{\linewidth}}}]{
                \texttt{\tiny{[GM$|$GM]}}
                \texttt{SPOSTA: S, (2,5)} \\
            }
        }
    }
    & & \\ \\

    \theutterance \stepcounter{utterance}  
    & & & \multicolumn{2}{p{0.3\linewidth}}{
        \cellcolor[rgb]{0.9,0.9,0.9}{
            \makecell[{{p{\linewidth}}}]{
                \texttt{\tiny{[GM$|$GM]}}
                \texttt{Penalità: (2,5) non è vuoto, ma contiene '┬'.} \\
\\ 
\texttt{You have collectively accumulated 15 of 16 penalties.} \\
\texttt{Riprova!} \\
            }
        }
    }
    & & \\ \\

    \theutterance \stepcounter{utterance}  
    & & \multicolumn{4}{p{0.6\linewidth}}{
        \cellcolor[rgb]{0.9,0.9,0.9}{
            \makecell[{{p{\linewidth}}}]{
                \texttt{\tiny{[P1$\langle$GM]}}
                \texttt{Penalità: (2,5) non è vuoto, ma contiene '┬'.} \\
\\ 
\texttt{You have collectively accumulated 15 of 16 penalties.} \\
\texttt{Riprova!} \\
            }
        }
    }
    & \\ \\

    \theutterance \stepcounter{utterance}  
    & \multicolumn{4}{p{0.6\linewidth}}{
        \cellcolor[rgb]{0.9,0.9,0.9}{
            \makecell[{{p{\linewidth}}}]{
                \texttt{\tiny{[P1$\rangle$GM]}}
                \texttt{SPOSTA: S, (3,5)} \\
            }
        }
    }
    & & \\ \\

    \theutterance \stepcounter{utterance}  
    & & & \multicolumn{2}{p{0.3\linewidth}}{
        \cellcolor[rgb]{0.9,0.9,0.9}{
            \makecell[{{p{\linewidth}}}]{
                \texttt{\tiny{[GM$|$GM]}}
                \texttt{SPOSTA: S, (3,5)} \\
            }
        }
    }
    & & \\ \\

    \theutterance \stepcounter{utterance}  
    & & & \multicolumn{2}{p{0.3\linewidth}}{
        \cellcolor[rgb]{0.9,0.9,0.9}{
            \makecell[{{p{\linewidth}}}]{
                \texttt{\tiny{[GM$|$GM]}}
                \texttt{Penalità: (3,5) non è vuoto, ma contiene '─'.} \\
\\ 
\texttt{You have collectively accumulated 16 of 16 penalties.} \\
\texttt{Riprova!} \\
            }
        }
    }
    & & \\ \\

    \theutterance \stepcounter{utterance}  
    & & \multicolumn{4}{p{0.6\linewidth}}{
        \cellcolor[rgb]{0.9,0.9,0.9}{
            \makecell[{{p{\linewidth}}}]{
                \texttt{\tiny{[P1$\langle$GM]}}
                \texttt{Penalità: (3,5) non è vuoto, ma contiene '─'.} \\
\\ 
\texttt{You have collectively accumulated 16 of 16 penalties.} \\
\texttt{Riprova!} \\
            }
        }
    }
    & \\ \\

    \theutterance \stepcounter{utterance}  
    & \multicolumn{4}{p{0.6\linewidth}}{
        \cellcolor[rgb]{0.9,0.9,0.9}{
            \makecell[{{p{\linewidth}}}]{
                \texttt{\tiny{[P1$\rangle$GM]}}
                \texttt{SPOSTA: O, (4,6)} \\
            }
        }
    }
    & & \\ \\

    \theutterance \stepcounter{utterance}  
    & & & \multicolumn{2}{p{0.3\linewidth}}{
        \cellcolor[rgb]{0.9,0.9,0.9}{
            \makecell[{{p{\linewidth}}}]{
                \texttt{\tiny{[GM$|$GM]}}
                \texttt{SPOSTA: O, (4,6)} \\
            }
        }
    }
    & & \\ \\

    \theutterance \stepcounter{utterance}  
    & & & \multicolumn{2}{p{0.3\linewidth}}{
        \cellcolor[rgb]{0.9,0.9,0.9}{
            \makecell[{{p{\linewidth}}}]{
                \texttt{\tiny{[GM$|$GM]}}
                \texttt{Penalità: (4,6) non è vuoto, ma contiene '│'.} \\
\\ 
\texttt{You have collectively accumulated 17 of 16 penalties.} \\
\texttt{Riprova!} \\
            }
        }
    }
    & & \\ \\

    \theutterance \stepcounter{utterance}  
    & & & \multicolumn{2}{p{0.3\linewidth}}{
        \cellcolor[rgb]{0.9,0.9,0.9}{
            \makecell[{{p{\linewidth}}}]{
                \texttt{\tiny{[GM$|$GM]}}
                \texttt{Maximum number of penalties exceeded} \\
            }
        }
    }
    & & \\ \\

    \theutterance \stepcounter{utterance}  
    & & & \multicolumn{2}{p{0.3\linewidth}}{
        \cellcolor[rgb]{0.9,0.9,0.9}{
            \makecell[{{p{\linewidth}}}]{
                \texttt{\tiny{[GM$|$GM]}}
                \texttt{* success: False} \\
\texttt{* lose: True} \\
\texttt{* aborted: True} \\
\texttt{{-}{-}{-}{-}{-}{-}{-}} \\
\texttt{* Shifts: 7.00} \\
\texttt{* Max Shifts: 12.00} \\
\texttt{* Min Shifts: 6.00} \\
\texttt{* End Distance Sum: 21.61} \\
\texttt{* Init Distance Sum: 30.19} \\
\texttt{* Expected Distance Sum: 29.33} \\
\texttt{* Penalties: 17.00} \\
\texttt{* Max Penalties: 16.00} \\
\texttt{* Rounds: 6.00} \\
\texttt{* Max Rounds: 28.00} \\
\texttt{* Object Count: 7.00} \\
            }
        }
    }
    & & \\ \\

    \theutterance \stepcounter{utterance}  
    & & & \multicolumn{2}{p{0.3\linewidth}}{
        \cellcolor[rgb]{0.9,0.9,0.9}{
            \makecell[{{p{\linewidth}}}]{
                \texttt{\tiny{[GM$|$GM]}}
                \texttt{Initial Grids:} \\
\texttt{```} \\
\texttt{Player 1:     Player 2:} \\
\texttt{ 1234567       1234567} \\
\texttt{╔═══╤═╤═╗     ╔═══╤═╤═╗} \\
\texttt{║◌◌◌│◌│◌║ 1   ║◌MS│◌│◌║ 1} \\
\texttt{║◌◌M│◌├─╢ 2   ║◌O◌│◌├─╢ 2} \\
\texttt{║AT◌├─┘◌║ 3   ║◌◌◌├─┘◌║ 3} \\
\texttt{║◌◌P│◌◌◌║ 4   ║◌◌◌│◌◌◌║ 4} \\
\texttt{╟─┬─┼───╢ 5   ╟─┬─┼───╢ 5} \\
\texttt{║◌│◌│S◌O║ 6   ║◌│◌│◌◌D║ 6} \\
\texttt{║◌│◌│◌D◌║ 7   ║T│◌│◌PA║ 7} \\
\texttt{╚═╧═╧═══╝     ╚═╧═╧═══╝} \\
\texttt{```} \\
\\ 
\texttt{Final Grids:} \\
\texttt{```} \\
\texttt{Player 1:     Player 2:} \\
\texttt{ 1234567       1234567} \\
\texttt{╔═══╤═╤═╗     ╔═══╤═╤═╗} \\
\texttt{║◌◌◌│◌│◌║ 1   ║◌◌◌│◌│◌║ 1} \\
\texttt{║◌◌◌│◌├─╢ 2   ║◌◌◌│◌├─╢ 2} \\
\texttt{║◌◌◌├─┘◌║ 3   ║◌◌◌├─┘◌║ 3} \\
\texttt{║A◌M│TP◌║ 4   ║◌◌S│O◌◌║ 4} \\
\texttt{╟─┬─┼───╢ 5   ╟─┬─┼───╢ 5} \\
\texttt{║◌│◌│S◌O║ 6   ║M│◌│T◌D║ 6} \\
\texttt{║◌│◌│◌D◌║ 7   ║◌│◌│◌PA║ 7} \\
\texttt{╚═╧═╧═══╝     ╚═╧═╧═══╝} \\
\texttt{```} \\
            }
        }
    }
    & & \\ \\

\end{supertabular}
}

\end{document}
